% Relatório Técnico da Semana 3 (issue #40). Fonte do texto: RELATORIO_TECNICO_SEMANA_03.md (mesma pasta).
% Compilar com pdflatex (Overleaf). Mesmo preâmbulo do relatório da Semana 2.
\documentclass[11pt,a4paper]{article}
\usepackage[utf8]{inputenc}
\usepackage[T1]{fontenc}
\usepackage[brazil]{babel}
\usepackage{mathptmx}
\usepackage[left=2.5cm,top=2.2cm,right=2.0cm,bottom=2.0cm]{geometry}
\usepackage{setspace}
\setstretch{1.08}
\usepackage{indentfirst}
\setlength{\parindent}{1.0cm}
\usepackage{titlesec}
\titleformat{\section}{\bfseries\normalsize}{\thesection}{1em}{\MakeUppercase}
\titleformat{\subsection}{\bfseries\normalsize}{\thesubsection}{1em}{}
\titlespacing*{\section}{0pt}{8pt}{3pt}
\titlespacing*{\subsection}{0pt}{6pt}{2pt}
\usepackage{fancyhdr}
\pagestyle{fancy}
\fancyhf{}
\fancyhead[R]{\thepage}
\renewcommand{\headrulewidth}{0pt}
\setlength{\headheight}{14pt}
\usepackage{caption}
\captionsetup[table]{name=Quadro,labelsep=endash,justification=raggedright,singlelinecheck=false,font=bf,skip=2pt}
\usepackage{booktabs}
\usepackage{tabularx}
\usepackage{array}
\usepackage{amsmath}
\usepackage{amssymb}
\usepackage{enumitem}
\setlist{nosep,leftmargin=0.7cm,itemsep=1pt,topsep=1pt,partopsep=0pt,parsep=0pt}
\usepackage{microtype}
\usepackage[table]{xcolor}
\definecolor{linkblue}{HTML}{0969DA}
\definecolor{highlightbg}{HTML}{F0F6FF}
\definecolor{highlightborder}{HTML}{1D4ED8}
\usepackage[colorlinks=true,linkcolor=black,citecolor=black,urlcolor=linkblue]{hyperref}
\urlstyle{same}
\sloppy
\newcolumntype{C}{>{\centering\arraybackslash}X}
\newcolumntype{L}{>{\raggedright\arraybackslash}X}
\newcolumntype{P}[1]{>{\centering\arraybackslash}p{#1}}
\newcommand{\fonte}{\vspace{0.1cm}\flushleft{\footnotesize{Fonte: elaborado pelos autores (2026).}}}
\newcommand{\codigo}[1]{\texttt{\small #1}}
\begin{document}

\pagenumbering{alph}
\begin{titlepage}
\begin{center}
{\bfseries INSTITUTO DE PESQUISAS ELDORADO\\
UNIVERSIDADE DE BRASÍLIA\\
RESIDÊNCIA EM INTELIGÊNCIA ARTIFICIAL -- BRASÍLIA\\[2.8cm]

GRUPO 07\\[0.35cm]
LETÍCIA MITI TAKAHASHI -- CIENTISTA DE DADOS\\
MARIA EDUARDA -- ML/LLMOPS \& DEVOPS\\
MARIA CAROLINA -- ENGENHARIA DE DADOS\\
CAUAN SOARES -- ARQUITETO DA SOLUÇÃO\\
SAMARA LETÍCIA -- ENGENHARIA DE DADOS\\[3.0cm]

ORQUESTRADOR DE ROLÊS (DISTRITO FEDERAL)\\[0.3cm]
{\normalsize RELATÓRIO TÉCNICO: COMPARAÇÃO ENTRE LLM PURA E RAG\\
Respostas baseadas em evidências -- Challenge 2, Semana 3}\\[1.2cm]

\fcolorbox{highlightborder}{highlightbg}{%
\parbox{0.92\textwidth}{\centering
\vspace{6pt}
{\normalsize\bfseries\textcolor{linkblue}{COMPARAÇÃO INTERATIVA: AS DUAS RESPOSTAS LADO A LADO}}\\[4pt]
{\large\bfseries\href{https://orquestrador-roles.onrender.com/comparativo}{\underline{https://orquestrador-roles.onrender.com/comparativo}}}\\[4pt]
{\footnotesize Central de Entregáveis: \href{https://orquestrador-roles.onrender.com/entregaveis}{https://orquestrador-roles.onrender.com/entregaveis}}
\vspace{6pt}
}}\\[2.8cm]

BRASÍLIA\\
2026}
\end{center}
\end{titlepage}

\pagenumbering{arabic}
\setcounter{page}{1}

\section{Resumo e desenho experimental}

O entregável da Semana 3 pede respostas baseadas em evidências: mostrar, com medição, se o RAG construído na Semana 2 responde melhor do que um modelo de linguagem sozinho. Testamos seis configurações nas 20 perguntas do \textit{golden dataset}, cada uma rodada duas vezes, com as respostas avaliadas por código e por um segundo modelo que atua como juiz.

\textbf{Resultado principal.} A LLM pura não acertou o lugar do gabarito em nenhuma das 18 perguntas sobre o DF (0 de 18). O RAG com a busca por palavras acertou 10 de 18, e com a busca por \textit{embeddings}, 16 de 18. Nas configurações com RAG, a fidelidade à base ficou entre 96,8\% e 100\%, acima da meta de 95\% do RNF-03.

\textbf{Achado que orienta a próxima semana.} Com a busca por palavras, colocar o LLM depois da busca não aumentou o acerto: o texto fixo da Semana 2 e o RAG com LLM acertaram os mesmos 12 de 20. Quem limita o acerto é a recuperação: o lugar certo só chega entre os 3 documentos em 10 de 18 perguntas. Com a busca por \textit{embeddings}, isso sobe para 16 de 18, mas a latência passa do limite de 5 s do RNF-14.

\textbf{Mesmo modelo nos dois modos.} O Modo A (LLM pura) e o Modo B (RAG) usam o mesmo modelo, o DeepSeek V4.1 Flash (\codigo{deepseek/deepseek-v4.1-flash-1017}, pela OpenRouter), com temperatura 0, o mesmo limite de palavras e a mesma instrução de não inventar. A única diferença é que o Modo B recebe até 3 lugares da base como \codigo{[FONTE \#n]} e precisa citá-los. Assim, a diferença entre os modos mede a base, e não o modelo nem o \textit{prompt} (ADR 0011).

\textbf{Por que não o Gemini 1.5 Flash.} O plano inicial previa o Gemini 1.5 Flash, que foi descontinuado e para o qual o grupo não tinha chave. O DeepSeek V4.1 Flash custa cerca de R\$ 0,0008 por pergunta, responde em 3 a 4 s e permite ligar o raciocínio, o que virou uma das configurações testadas (ADR 0010). As chamadas usam o roteamento sem retenção de dados da OpenRouter (\codigo{zdr}, \codigo{data\_collection: deny}).

\section{Golden dataset}

O \textit{golden dataset} (\codigo{data/golden\_dataset\_semana3.json}) tem 20 perguntas escritas como um usuário faria, cada uma com um gabarito copiado da base de 93 lugares validados em dois fatores: nome, região administrativa (RA), preço médio e horário, além do vídeo do TikTok que serve de evidência. O Quadro~\ref{tab:golden} mostra as categorias.

\begin{table}[!htbp]
\centering
\small
\renewcommand{\arraystretch}{1.03}
\caption{Categorias do golden dataset}
\label{tab:golden}
\begin{tabularx}{\textwidth}{p{5.0cm} c L}
\toprule
\textbf{Categoria} & \textbf{Casos} & \textbf{O que testa} \\
\midrule
RA fora do centro & 5 & Lugares reais fora do Plano Piloto (RNF-13) \\
Janela de horário & 3 & Horário real (madrugada depois da 1h, domingo) \\
Fora do DF & 2 & Recusar perguntas sobre outras cidades \\
Precisão financeira & 5 & Preço exato e teto estrito em reais (RNF-03) \\
Ambiente e restrições & 3 & Ambiente e restrição alimentar descritos só no vídeo \\
Conflito e gírias & 2 & Pedido que parece impossível e pergunta em candanguês (RNF-08) \\
\bottomrule
\end{tabularx}
\fonte
\end{table}

\begin{itemize}
\item \textbf{Uma resposta certa por pergunta.} Cada pergunta é específica o bastante para ter um único lugar certo na base: são 18 lugares diferentes para as 18 perguntas do DF.
\item \textbf{Gabarito travado por teste.} \codigo{tests/test\_golden\_dataset.py}, que roda no CI, confere que nome, RA e preço do gabarito são iguais aos da base e que as 6 categorias têm a distribuição acima.
\item \textbf{Armadilha registrada.} Cada caso anota o erro que esperamos da LLM pura, por exemplo ``sugerir hamburguerias do Plano Piloto'' para uma pergunta sobre Taguatinga.
\end{itemize}

\section{Configurações testadas}

\begin{table}[!htbp]
\centering
\small
\renewcommand{\arraystretch}{1.03}
\caption{As seis configurações do benchmark}
\label{tab:configs}
\begin{tabularx}{\textwidth}{c p{5.6cm} L}
\toprule
\textbf{\#} & \textbf{Configuração} & \textbf{Como funciona} \\
\midrule
1 & LLM pura, sem raciocínio & Modo A: o modelo responde só com o que sabe \\
2 & LLM pura, raciocínio alto & Modo A com \codigo{reasoning.effort = "high"} \\
3 & RAG sem LLM & Busca por palavras e o texto fixo da Semana 2, preenchido com o 1\textsuperscript{o} lugar \\
4 & RAG + LLM, busca por palavras & Modo B: busca por palavras e o modelo redige a resposta com as fontes \\
5 & RAG + LLM, palavras, raciocínio alto & Configuração 4 com o raciocínio alto \\
6 & RAG + LLM, busca por \textit{embeddings} & Configuração 4 com a busca vetorial (\codigo{google/gemini-embedding-001}) \\
\bottomrule
\end{tabularx}
\fonte
\end{table}

As três configurações com LLM e RAG compartilham os mesmos ajustes da busca e os mesmos \textit{guardrails}, que vêm antes do modelo e são escritos em Python:

\begin{itemize}
\item \textbf{3 documentos por pergunta.} O modelo recebe no máximo 3 lugares, com nome, tipo, RA, endereço, horário, faixa e preço médio e descrição.
\item \textbf{Limiar de aderência.} Se a maior pontuação da busca fica abaixo do limiar, nenhum lugar é recomendado e o modelo nem é chamado. O limiar é 0,65 na busca por palavras e 0,63 no cosseno da busca vetorial. Os dois foram calibrados com perguntas fora do tema da base. Na busca por palavras, essas perguntas ficaram entre 0,50 e 0,56. No cosseno, ficaram entre 0,574 e 0,624, e as perguntas do \textit{golden}, entre 0,625 e 0,812 (ADR 0012).
\item \textbf{Escopo do DF.} Perguntas que citam um lugar fora do DF, sem nenhuma RA, são recusadas antes da busca.
\item \textbf{Teto em reais.} Quando a pergunta tem teto ou pede algo grátis, os lugares acima do teto são retirados em Python, antes de chegar ao modelo.
\end{itemize}

A busca híbrida (palavras e \textit{embeddings} juntos, por \textit{Reciprocal Rank Fusion}) foi medida na prévia de 12/11 e ficou de fora do benchmark, porque não superou a vetorial no acerto (6 de 8 contra 7 de 8). Como teste de ferramenta, também montamos a busca por palavras no framework Haystack (BM25): ela traz o lugar certo entre os 3 primeiros em 14 de 18 perguntas. É o mesmo tipo de busca, então o resultado não diz nada sobre o framework em si.

\section{Como as respostas foram avaliadas}

O avaliador (\codigo{src/telemetria/evals\_semana3.py}) funciona em duas etapas.

\begin{enumerate}
\item \textbf{Verificação em código (RNF-04).} Antes do juiz, o código confere se os lugares citados existem entre os 93 da base, se a RA do lugar recomendado é a esperada, se os preços citados batem com os da base e se a pergunta fora do DF foi recusada.
\item \textbf{Juiz.} O MiMo-V2.6-Pro (\codigo{xiaomi/mimo-v2.6-pro}), de outra família que o modelo avaliado, roda com temperatura 0 e saída em JSON. Ele classifica cada afirmação da resposta sobre lugar, preço, endereço ou horário como ``confere com a base'', ``contradiz a base'' ou ``não está na base''. Para o Modo B, o juiz recebe os mesmos documentos que o modelo recebeu; para o Modo A, os lugares citados que existem na base e os do gabarito. Usamos um juiz de outra família porque um modelo tende a favorecer as próprias respostas.
\end{enumerate}

\begin{table}[!htbp]
\centering
\small
\renewcommand{\arraystretch}{1.03}
\caption{Notas calculadas para cada resposta}
\label{tab:notas}
\begin{tabularx}{\textwidth}{p{4.4cm} L}
\toprule
\textbf{Nota} & \textbf{O que mede} \\
\midrule
Acerto do gabarito & O lugar do gabarito aparece na resposta (fora do DF, a recusa conta como acerto) \\
Fidelidade (RNF-03) & Proporção das afirmações que conferem com a base \\
Alucinação & Proporção das afirmações que a base desmente \\
Fora da base & Proporção das afirmações que a base não permite verificar \\
Precisão da RA (RNF-13) & Algum lugar recomendado está na RA esperada \\
Recall@3 e MRR & O lugar do gabarito está entre os 3 documentos recuperados, e em que posição \\
Latência (RNF-14) & Tempo de resposta, com meta de p90 até 5 s \\
Custo (RNF-15) & Custo das chamadas pagas por pergunta, em reais (cotação de R\$ 5,80 por dólar) \\
\bottomrule
\end{tabularx}
\fonte
\end{table}

Uma recusa ou um ``não sei'' não tem afirmação a medir. Nesse caso, a resposta conta como abstenção nos dois modos, fica sem nota de fidelidade e de alucinação, e o erro, quando havia resposta na base, aparece no acerto do gabarito.

\section{Resultados}

\subsection{Comparação das seis configurações}

O Quadro~\ref{tab:resultados} traz a rodada 1, com 20 perguntas por configuração. As contagens são sobre os casos em que a nota se aplica. O p90 é calculado sobre as 20 latências com \codigo{statistics.quantiles} do Python (método exclusivo, que interpola entre a 18\textsuperscript{a} e a 19\textsuperscript{a}).

\begin{table}[!htbp]
\centering
\scriptsize
\setlength{\tabcolsep}{3pt}
\renewcommand{\arraystretch}{1.1}
\caption{Resultados das seis configurações (rodada 1)}
\label{tab:resultados}
\begin{tabular}{c >{\raggedright\arraybackslash}p{3.1cm} P{1.15cm} P{1.15cm} P{1.15cm} P{0.95cm} P{1.1cm} P{1.25cm} P{0.7cm} P{1.55cm} P{1.0cm}}
\toprule
\textbf{\#} & \textbf{Configuração} & \textbf{Acerto no DF (18)} & \textbf{Fideli\-dade} & \textbf{Aluci\-nação} & \textbf{Fora da base} & \textbf{Precisão da RA} & \textbf{Recall@3} & \textbf{MRR} & \textbf{Latência p50 / p90 (s)} & \textbf{Custo (R\$)} \\
\midrule
1 & LLM pura & 0 de 18 & 0,0\% & 3,0\% & 97,0\% & 0 de 15 & -- & -- & 3,75 / 5,26 & 0,00079 \\
2 & LLM pura, raciocínio alto & 0 de 18 & 0,0\% & 28,6\%\textsuperscript{1} & 71,4\% & 0 de 14 & -- & -- & 6,84 / 14,58 & 0,00198 \\
3 & RAG sem LLM & 10 de 18 & 99,2\% & 0,0\% & 0,8\% & 13 de 16 & 10 de 18 & 0,37 & 0,02 / 0,03 & 0 \\
4 & RAG + LLM, palavras & 10 de 18 & 98,7\% & 0,0\% & 1,4\% & 13 de 16 & 10 de 18 & 0,37 & 3,15 / 3,61 & 0,00081 \\
5 & RAG + LLM, palavras, raciocínio alto & 10 de 18 & 100,0\% & 0,0\% & 0,0\% & 13 de 16 & 10 de 18 & 0,37 & 4,51 / 6,88 & 0,00121 \\
6 & RAG + LLM, \textit{embeddings} & 16 de 18 & 96,8\% & 0,7\% & 2,5\% & 16 de 17 & 16 de 18 & 0,83 & 5,10 / 5,83 & 0,00092 \\
\bottomrule
\end{tabular}
\vspace{0.1cm}
\flushleft{\footnotesize{Nas 2 perguntas fora do DF, todas as configurações recusaram, o que conta como acerto: o acerto nas 20 perguntas é o do DF mais 2 (por exemplo, 12 de 20 na configuração 4).\\
\textsuperscript{1} Valor inflado por um defeito do avaliador (Seção~\ref{sec:limitacoes}): as quatro respostas que somam essa taxa são recusas que só repetem o teto da pergunta.\\
Fonte: elaborado pelos autores (2026), a partir de \codigo{data/resultados\_benchmark\_semana3.json}.}}
\end{table}

\begin{itemize}
\item \textbf{A LLM pura não erra o preço: ela fala de lugares que não pode comprovar.} Na configuração 1, 97\% das afirmações citam lugares fora da base. São 32 nomes em 13 das 20 respostas: parte existe, como o Coco Bambu; parte são descrições genéricas (``Feira de Vicente Pires'', ``praça de alimentação de shopping próximo'') ou palavras que o extrator tomou por nome (``Alternativa''). Contradições reais, em que a base desmente o que foi dito, apareceram em 3 respostas. Com o raciocínio alto, o modelo se recusa em 6 das 20 perguntas (5 sem o raciocínio) e continua sem acertar nenhum lugar do DF.
\item \textbf{O RAG fica fiel à base.} Nas configurações 3 a 6, nenhum lugar fora da base foi recomendado. A única contradição da rodada 1 está na configuração 6 e é discutida na Seção~\ref{sec:casos}.
\item \textbf{O raciocínio alto não ajuda o RAG.} A configuração 5 acerta o mesmo que a 4, custa 50\% mais e passa dos 5 s em 7 de 20 perguntas.
\item \textbf{A busca por \textit{embeddings} acerta mais, mas é mais lenta.} O \textit{embedding} da pergunta leva cerca de 1,5 s pela rede, e 11 de 20 respostas passam de 5 s.
\end{itemize}

\subsection{Resultados por categoria}

\begin{table}[!htbp]
\centering
\small
\renewcommand{\arraystretch}{1.03}
\caption{Acerto do gabarito por categoria (rodada 1)}
\label{tab:categorias}
\begin{tabularx}{\textwidth}{p{3.8cm} c C C C C}
\toprule
\textbf{Categoria} & \textbf{Casos} & \textbf{LLM pura (1)} & \textbf{RAG sem LLM (3)} & \textbf{RAG + LLM, palavras (4)} & \textbf{RAG + LLM, \textit{embeddings} (6)} \\
\midrule
RA fora do centro & 5 & 0 de 5 & 4 de 5 & 4 de 5 & 5 de 5 \\
Janela de horário & 3 & 0 de 3 & 2 de 3 & 2 de 3 & 2 de 3 \\
Fora do DF & 2 & 2 de 2 & 2 de 2 & 2 de 2 & 2 de 2 \\
Precisão financeira & 5 & 0 de 5 & 3 de 5 & 3 de 5 & 5 de 5 \\
Ambiente e restrições & 3 & 0 de 3 & 1 de 3 & 1 de 3 & 3 de 3 \\
Conflito e gírias & 2 & 0 de 2 & 0 de 2 & 0 de 2 & 1 de 2 \\
\bottomrule
\end{tabularx}
\fonte
\end{table}

A busca por \textit{embeddings} ganha sobretudo em preço, ambiente e restrição alimentar, perguntas que descrevem a experiência mais do que o nome ou o tipo do lugar. Nessas categorias, 4 perguntas só foram acertadas pela busca vetorial: ``jantar romântico a dois num restaurante chique'' (Lucca), ``rodízio de churrasco japonês'' (Yaki Nikuya), ``doceria sem trigo, sem lactose e sem açúcar refinado'' (Lefit) e ``prato executivo em Samambaia por até R\$ 30'' (Saguari).

\subsection{Variação entre as duas rodadas}

As duas rodadas foram feitas sem cache.

\begin{itemize}
\item \textbf{Acerto estável.} O acerto do gabarito foi idêntico nas duas rodadas, em todas as configurações.
\item \textbf{Variação pequena no RAG.} Nas configurações 4 a 6, a nota do juiz mudou em no máximo 3 das 20 perguntas, e a alucinação média variou menos de 1 ponto percentual.
\item \textbf{Variação maior na LLM pura.} A latência da LLM pura variou mais: o p90 da configuração 2 foi de 14,6 s para 18,1 s.
\end{itemize}

\section{Critério de escolha e configuração escolhida}

O critério foi definido antes de rodar o benchmark (issue \#38): entre as configurações com p90 de até 5 s e nenhum lugar recomendado fora da base, vence a de maior acerto do gabarito, e no empate, a mais barata.

\begin{itemize}
\item \textbf{Configurações elegíveis.} Só a 3 e a 4 passam pelos dois filtros. A 1 e a 2 recomendam lugares fora da base; a 5 e a 6 passam de 5 s no p90.
\item \textbf{Vencedora pelo critério: configuração 3, RAG sem LLM.} Ela empata com a 4 em 12 de 20 de acerto e vence no desempate por não ter custo de modelo.
\item \textbf{Configuração adotada no produto: a 4, RAG + LLM com a busca por palavras.} Fica dentro do RNF-14 (p90 de 3,61 s), tem 98,7\% de fidelidade e custa R\$ 0,0008 por pergunta. É uma decisão de produto, à parte do critério. O texto fixo só sabe recomendar o 1\textsuperscript{o} lugar da busca; o LLM explica quando o lugar atende só em parte (``isso ultrapassa seus R\$ 30 em cerca de R\$ 3''), contextualiza as recusas e escolhe entre os 3 documentos o que de fato responde à pergunta.
\item \textbf{Próximo passo: a busca vetorial.} Ela é a candidata para depois que o \textit{embedding} da pergunta ficar mais rápido, com cache ou um modelo local. Com 16 de 18 de \textit{recall}, é a configuração que mais acerta. O ganho estimado no p90 ainda não foi medido.
\end{itemize}

\section{Análise qualitativa de casos}
\label{sec:casos}

\subsection{Caso 1 -- RA fora do centro (golden\_01)}

\textit{``Onde tem rodízio de mini hambúrguer e pão de queijo em Taguatinga para ir com amigos?''}

\begin{itemize}
\item \textbf{LLM pura.} Disse que não tinha certeza e, mesmo assim, sugeriu ``Bar do Caldo (Taguatinga)'', ``Boteco da Asa (Taguatinga)'' e o Taguatinga Shopping, sempre com ``não sei confirmar''. Nenhum está na base.
\item \textbf{RAG + LLM.} Indicou o Filomenambúrguer [FONTE \#2] em Taguatinga Norte, a R\$ 49,99, com o rodízio de pão de queijo e a brinquedoteca descritos no \href{https://www.tiktok.com/@rolezandoembrasilia/video/7626405885388197128}{vídeo de evidência}. O lugar do gabarito era o 2\textsuperscript{o} da busca, e o modelo o escolheu em vez do 1\textsuperscript{o}.
\end{itemize}

\subsection{Caso 2 -- Quando a busca falha (golden\_05)}

\textit{``Onde comer rodízio de pizza e sushi à noite em Samambaia com a família?''}

\begin{itemize}
\item \textbf{Busca por palavras.} Não trouxe o VBQ. O modelo respondeu, com honestidade, que a base não tinha um lugar com rodízio de pizza e sushi em Samambaia, e listou as opções mais próximas em outras RAs.
\item \textbf{Busca por \textit{embeddings}.} Trouxe o VBQ em 1\textsuperscript{o} lugar (Samambaia, R\$ 44,90, mais de 30 sabores de pizza e ilha de sushi, \href{https://www.tiktok.com/@rolezandoembrasilia/video/7678112514009304338}{vídeo}), e a resposta acertou.
\item \textbf{Leitura.} O erro é da recuperação, não da geração.
\end{itemize}

\subsection{Caso 3 -- Recusa indevida (golden\_06)}

\textit{``Onde dá para comer e beber com os amigos depois da 1h da manhã num sábado?''}

\begin{itemize}
\item \textbf{O que aconteceu.} As duas buscas recusaram a pergunta: a aderência ficou em 0,62 na busca por palavras (limiar de 0,65) e em 0,6248 no cosseno (limiar de 0,63). O New Mercadito, aberto até 1h30, está na base.
\item \textbf{Leitura.} O limiar que protege contra perguntas fora do tema também barra uma pergunta válida. A margem no cosseno é de cinco milésimos.
\end{itemize}

\subsection{Caso 4 -- Fora do DF (golden\_09)}

\textit{``Quais são os melhores bares em Copacabana, no Rio de Janeiro, para ir com amigos numa sexta à noite?''}

\begin{itemize}
\item \textbf{RAG.} O \textit{guardrail} recusou em menos de 1 ms, sem chamar o modelo.
\item \textbf{LLM pura.} Também recusou, porque o \textit{prompt} a apresenta como assistente de rolês no DF, e se ofereceu para sugerir bares em Brasília.
\item \textbf{Leitura.} As duas recusas contam como acerto.
\end{itemize}

\subsection{Caso 5 -- A única contradição do RAG (golden\_15)}

\textit{``Tenho só R\$ 30, nem um centavo a mais. Dá para almoçar no buffet à vontade com tempero nordestino em Vicente Pires numa terça-feira?''}

\begin{itemize}
\item \textbf{O que o modelo respondeu.} Com a busca vetorial, ``Sim, dá.'' e, na frase seguinte, que o preço médio do Restaurante do Daniel é R\$ 32,99, acima dos R\$ 30.
\item \textbf{Por que o juiz marcou contradição.} O ``Sim, dá'' pressupõe um preço de até R\$ 30, e o documento traz R\$ 32,99.
\item \textbf{Leitura.} Nenhum dado foi inventado, e a própria resposta corrige o preço na frase seguinte. O problema é a resposta começar por uma conclusão errada. Com a busca por palavras, o mesmo modelo começou com ``Sim, dá -- mas por pouco'' e foi marcado como fiel. A diferença entre as duas frases é pequena, e é nesse tipo de caso que o juiz e uma leitura humana podem discordar.
\end{itemize}

\subsection{Caso 6 -- Guardrail literal demais (golden\_19)}

\textit{``Quero comemorar meu aniversário num rodízio japonês completo na Asa Norte, mas sem pagar nada. Existe rodízio caro totalmente de graça?''}

\begin{itemize}
\item \textbf{O que o RAG respondeu.} Leu ``sem pagar nada'' como teto de R\$ 0, retirou todos os lugares e respondeu que nenhum lugar da base é grátis.
\item \textbf{O que a base tem.} O gabarito é o Blue Sushi: segundo o vídeo, o aniversariante que leva 10 convidados não paga o próprio rodízio.
\item \textbf{Leitura.} O filtro de teto, que protege o orçamento, não sabe ler uma condição.
\end{itemize}

\section{Limitações}
\label{sec:limitacoes}

\begin{itemize}
\item \textbf{Amostra pequena.} São 20 perguntas e uma base de 93 lugares. Uma pergunta vale 5 pontos percentuais de acerto, então diferenças de uma ou duas perguntas não devem ser lidas como tendência.
\item \textbf{Defeito no avaliador da LLM pura com raciocínio.} Os 28,6\% de alucinação da configuração 2 vêm de quatro recusas (golden\_07, 14, 15 e 20) que só repetem o teto da pergunta, como em ``não tenho informação confiável para indicar buffet\ldots{} por menos de R\$ 30''. O código toma esse valor como preço inventado. Lidas à mão, essas respostas não contêm contradição.
\item \textbf{Juiz diferente na rodada 1 da LLM pura.} Na rodada 1, as configurações 1 e 2 foram julgadas com a versão anterior do \textit{prompt} do juiz, e as configurações 3 a 6, com a versão corrigida. A rodada 2 usa a versão corrigida em todas.
\item \textbf{``Fora da base'' não é invenção comprovada.} A nota mostra o que não conseguimos verificar. Parte dos lugares citados pela LLM pura existe.
\item \textbf{O juiz também erra.} Antes da correção do \textit{prompt}, ele marcava como contradição respostas que diziam ``a base não tem'', comparando-as com documentos que o modelo não tinha recebido. As contradições que sobraram foram conferidas à mão (Seção~\ref{sec:casos}).
\item \textbf{Uma data e um modelo.} Os resultados valem para o DeepSeek V4.1 Flash, as versões de \textit{prompt} de 10/11 (Modo A) e 12/11 (Modo B) e a base de novembro de 2026.
\end{itemize}

\section{Telemetria e auditoria}

\begin{itemize}
\item \textbf{Langfuse.} Cada comparação gera um \textit{trace} no Langfuse Cloud, com os dois modos, a busca, os \textit{guardrails} e as notas do avaliador. O benchmark foi registrado como seis experimentos sobre o \textit{dataset} \codigo{golden\_dataset\_semana3}, um por configuração, com o \textit{trace} de cada resposta. Os links estão no log do benchmark (\codigo{BENCHMARK\_LLM\_VS\_RAG\_LOG.txt}); o projeto no Langfuse é de acesso do grupo.
\item \textbf{Registro sem dados pessoais (RNF-10).} O \codigo{logger\_metricas} grava em JSONL a latência, os \textit{tokens}, o custo e o modo de cada chamada, e mascara CPF, e-mail, telefone e valores de orçamento.
\item \textbf{Cache (RNF-15).} O cache por SHA-256 da pergunta evita pagar duas vezes pela mesma chamada; o benchmark foi rodado com ele desligado.
\item \textbf{Reprodução.} \codigo{PYTHONPATH=. python scripts/executar\_benchmark\_semana3.py} refaz o benchmark. Os resultados estão em \codigo{data/resultados\_benchmark\_semana3.json} e \codigo{.csv}, e o log comprobatório, com as 20 respostas lado a lado, em \codigo{docs/entregas/semana-03/BENCHMARK\_LLM\_VS\_RAG\_LOG.txt}.
\item \textbf{Testes automáticos.} O \textit{golden dataset}, o avaliador, o comparador e a integridade dos resultados têm testes que rodam no CI do GitHub Actions.
\end{itemize}

\begin{center}
\fcolorbox{highlightborder}{highlightbg}{%
\parbox{0.96\textwidth}{\centering
\vspace{4pt}
{\bfseries\normalsize\textcolor{linkblue}{TESTE A COMPARAÇÃO COM SUAS PRÓPRIAS PERGUNTAS}}\\[3pt]
\href{https://orquestrador-roles.onrender.com/comparativo}{\large\bfseries\underline{https://orquestrador-roles.onrender.com/comparativo}}\\[3pt]
{\footnotesize Central de Entregáveis: \href{https://orquestrador-roles.onrender.com/entregaveis}{https://orquestrador-roles.onrender.com/entregaveis}}
\vspace{4pt}
}}
\end{center}

\clearpage

\section*{Referências}
\addcontentsline{toc}{section}{Referências}

\noindent BRASIL. \textbf{Lei n\textsuperscript{o} 13.709, de 14 de agosto de 2018}. Lei Geral de Proteção de Dados Pessoais (LGPD). Diário Oficial da União: seção 1, Brasília, DF, 15 ago. 2018.

\vspace{0.2cm}
\noindent CORMACK, Gordon V.; CLARKE, Charles L. A.; BUETTCHER, Stefan. Reciprocal rank fusion outperforms Condorcet and individual rank learning methods. In: \textbf{Proceedings of the 32nd International ACM SIGIR Conference on Research and Development in Information Retrieval}. 2009. p. 758-759.

\vspace{0.2cm}
\noindent LEWIS, Patrick et al. Retrieval-augmented generation for knowledge-intensive NLP tasks. \textbf{Advances in Neural Information Processing Systems}, v. 33, p. 9459-9474, 2020.

\vspace{0.2cm}
\noindent ZHENG, Lianmin et al. Judging LLM-as-a-judge with MT-Bench and Chatbot Arena. \textbf{Advances in Neural Information Processing Systems}, v. 36, 2023.

\end{document}
